\documentclass[11pt]{article}

\usepackage[T1]{fontenc}
\usepackage[utf8]{inputenc}
\usepackage{lmodern}
\usepackage{microtype}
\usepackage[margin=1.08in]{geometry}
\usepackage{amsmath,amssymb,amsthm,mathtools}
\usepackage{enumitem}
\usepackage{xcolor}
\usepackage{hyperref}

\hypersetup{
  colorlinks=true,
  linkcolor=blue!55!black,
  citecolor=blue!55!black,
  urlcolor=blue!60!black,
  pdftitle={A width-plateau counterexample to an index-p conjecture for mean-curvature flow}
}

\setlength{\parindent}{0pt}
\setlength{\parskip}{0.48em}
\setlength{\emergencystretch}{2em}
\setlist[itemize]{leftmargin=*,itemsep=0.25em,topsep=0.35em}

\newtheorem{theorem}{Theorem}
\newtheorem{proposition}[theorem]{Proposition}
\newtheorem{lemma}[theorem]{Lemma}
\newtheorem{corollary}[theorem]{Corollary}
\theoremstyle{remark}
\newtheorem{remark}[theorem]{Remark}

\newcommand{\Sph}{\mathbb S}
\newcommand{\RP}{\mathbb {RP}}
\newcommand{\Ztwo}{\mathbb Z_2}
\newcommand{\ind}{\operatorname{ind}}
\newcommand{\nul}{\operatorname{nul}}
\newcommand{\Ric}{\operatorname{Ric}}
\newcommand{\Area}{\operatorname{Area}}
\newcommand{\Mass}{\mathbf M}

\title{A Width-Plateau Counterexample to an\\
Index-$p$ Conjecture for Mean-Curvature Flow}
\author{[Author Name]}
\date{}

\begin{document}

\maketitle

\begin{abstract}
\textbf{Relation to the problem list.}
This manuscript addresses Question~55 (M5), ``An eternal flow leaves every
higher $p$-width hypersurface.''  It gives a full counterexample to the
literal positive-Ricci alternative in the published conjecture and in the
list.  It does not disprove a repaired version restricted to generic bumpy
metrics, nor does it settle the connecting-flow problem once a suitable
higher-index width realizer and lower-area endpoint are assumed to exist.
This file duplicates \texttt{round\_sphere\_pwidth\_counterexample.tex} and
is not a separate solution.

Chen and Gaspar conjectured that if a closed Riemannian manifold has a
bumpy metric or a metric of positive Ricci curvature, then, for every
$p\geq 2$, its $p$-width is realized by a min--max minimal hypersurface of
Morse index exactly $p$ which is connected by an eternal weak
mean-curvature flow to a minimal hypersurface of smaller area and smaller
index.  We observe that the positive-Ricci branch of the conjecture is
false.  On the unit round three-sphere one has
$\omega_2(\Sph^3)=4\pi$, and every Almgren--Pitts width realizer of
mass $4\pi$ is a
multiplicity-one equatorial two-sphere.  Its Morse index is $1$, not $2$.
Moreover, the round three-sphere has no nonempty closed minimal surface of
area below $4\pi$, so the required lower-area endpoint does not exist.
Thus the round three-sphere gives two independent obstructions to the
conjecture as published.
\end{abstract}

\section{Introduction}

Let $\omega_p(M,g)$ denote the $p$-width of a closed Riemannian manifold,
defined using mod-$2$ codimension-one cycles; see
\cite{MarquesNevesPositiveRicci,Pitts}.  Motivated by Morse theory for the
area functional, Chen and Gaspar formulated the following conjecture in
\cite[Conjecture~1]{ChenGaspar}.

\medskip
\noindent
\emph{Published conjecture.}
If $(M,g)$ is closed and $g$ is bumpy or has positive Ricci curvature,
then, for every integer $p\geq2$, there exists a min--max minimal
hypersurface $\Sigma_-$ such that
\[
  \Area(\Sigma_-)=\omega_p(M,g),
  \qquad \ind(\Sigma_-)=p,
\]
and $\Sigma_-$ can be connected by an eternal weak mean-curvature flow to
a minimal hypersurface $\Sigma_+$ having strictly smaller area and
strictly smaller Morse index.

The disjunction in the hypothesis is important: a positive-Ricci metric
need not be bumpy.  Symmetric metrics may have plateaux in their volume
spectrum, and on such a plateau the parameter $p$ need not equal the
Morse index of the realizing critical point.  The simplest example
already disproves the conjecture.

\begin{theorem}\label{thm:main}
Conjecture~1 of Chen--Gaspar \cite{ChenGaspar}, in its published form, is
false.  A counterexample is the unit round three-sphere with $p=2$.
Indeed, $\Ric_{\Sph^3}=2g_{\mathrm{rd}}>0$ and
$\omega_2(\Sph^3,g_{\mathrm{rd}})=4\pi$, but every minimal hypersurface
realizing this width is an equatorial $\Sph^2$, with
\[
  \ind(\Sph^2)=1,
  \qquad
  \nul(\Sph^2)=3.
\]
There is also no nonempty closed minimal surface in $\Sph^3$ of area
strictly less than $4\pi$.  Thus both the exact-index requirement and the
lower-area endpoint requirement fail.
\end{theorem}

Chen and Gaspar themselves record in the proof of
\cite[Proposition~1]{ChenGaspar} that the equatorial spheres realize the
first four widths of the round $\Sph^3$.  We nevertheless give a direct
argument for the only identity needed here, $\omega_2=4\pi$, so that the
counterexample is self-contained modulo standard Almgren--Pitts
existence and regularity.

\section{The second width of the round three-sphere}

Write
\[
  \mathcal Z_2=\mathcal Z_2(\Sph^3;\Ztwo)
\]
for the space of flat mod-$2$ two-cycles in $\Sph^3$, and let
$\bar\lambda\in H^1(\mathcal Z_2;\Ztwo)$ be its canonical generator.  A
continuous map $\Phi:X\to\mathcal Z_2$ is a $p$-sweepout when
\[
  \Phi^*(\bar\lambda^p)\neq0
\]
and the usual no-concentration condition holds.

\begin{lemma}\label{lem:width}
For the unit round three-sphere,
\[
  \omega_1(\Sph^3)=\omega_2(\Sph^3)=4\pi.
\]
\end{lemma}

\begin{proof}
The standard family of geodesic spheres, starting and ending at the zero
cycle, is a one-sweepout whose largest member is an equator.  Hence
$\omega_1(\Sph^3)\leq4\pi$.

Conversely, Almgren--Pitts min--max theory \cite{Pitts} produces a nonzero
stationary integral varifold $V$ with mass
$\|V\|(\Sph^3)=\omega_1(\Sph^3)$.  In ambient dimension three its support
is a smooth embedded closed minimal surface, possibly with positive
integer multiplicities.  By stereographic projection and conformal
invariance, the classical Euclidean Willmore inequality becomes the
following spherical inequality \cite{Willmore,MarquesNevesWillmore}: every
closed immersed surface $\Sigma\subset\Sph^3$ satisfies
\begin{equation}\label{eq:willmore}
  \Area(\Sigma)+\frac14\int_\Sigma |\mathbf H|^2\,d\mu\geq4\pi,
\end{equation}
where $\mathbf H$ is the trace mean-curvature vector; equality occurs only
for a geodesic two-sphere.  Each minimal component of $V$ therefore has
area at least $4\pi$.  It follows that $\omega_1(\Sph^3)\geq4\pi$, and
thus $\omega_1(\Sph^3)=4\pi$.

It remains to give a two-sweepout of maximal mass $4\pi$.  Regard
$\RP^2$ as a projective plane in the $\RP^3$ of normal lines in
$\mathbb R^4$, and define
\[
  \Phi:\RP^2\longrightarrow\mathcal Z_2,
  \qquad
  \Phi([v])=\Sph^3\cap v^\perp.
\]
Every member is an equatorial two-sphere and hence has mass $4\pi$.  If
$\alpha$ denotes the generator of $H^1(\RP^2;\Ztwo)$, then
\begin{equation}\label{eq:pullback}
  \Phi^*\bar\lambda=\alpha.
\end{equation}
The map $\Phi$ is continuous in the flat topology.  It also has no
concentration of mass: for some universal $C$ and all sufficiently small
$r>0$,
\[
  \sup_{[v]\in\RP^2}\sup_{x\in\Sph^3}
  \|\Phi([v])\|(B_r(x))\leq Cr^2\longrightarrow0.
\]
After choosing any finite cubulation of $\RP^2$, it is therefore an
admissible family in the definition of the width.

To verify \eqref{eq:pullback}, restrict $\Phi$ to a projective line.  The
resulting loop rotates an equator through a half turn.  Equivalently, a
generic point of $\Sph^3$ belongs to exactly one equator in the loop,
counted modulo $2$.  Its Almgren swept chain is consequently the
fundamental class $[\Sph^3]$, so Almgren's isomorphism evaluates
$\bar\lambda$ nontrivially on this loop.  Hence
$\Phi^*\bar\lambda=\alpha$.  Therefore
\[
  \Phi^*(\bar\lambda^2)=\alpha^2\neq0
  \quad\text{in }H^2(\RP^2;\Ztwo).
\]
Thus $\Phi$ is a two-sweepout, and
$\omega_2(\Sph^3)\leq4\pi$.  Since the widths are nondecreasing,
\[
  4\pi=\omega_1(\Sph^3)\leq\omega_2(\Sph^3)\leq4\pi,
\]
which proves the claim.
\end{proof}

\begin{remark}
The same phenomenon is usually stated as
\[
  \omega_1(\Sph^3)=\omega_2(\Sph^3)
  =\omega_3(\Sph^3)=\omega_4(\Sph^3)=4\pi,
  \qquad
  \omega_5(\Sph^3)=2\pi^2;
\]
see, for example, \cite{Nurser} and the discussion in
\cite{ChenGaspar}.  Only the first equality is required below.
\end{remark}

\section{Rigidity of the width realizers}

The mass value $4\pi$ leaves no alternative realizing surface.

\begin{lemma}\label{lem:rigidity}
Let $V$ be a stationary integral two-varifold in the unit round
$\Sph^3$ whose support has the regularity of an Almgren--Pitts min--max
hypersurface.  If $\|V\|(\Sph^3)=4\pi$, then
\[
  V=|E|
\]
for an equatorial two-sphere $E$, with multiplicity one.
\end{lemma}

\begin{proof}
By regularity, write
\[
  V=\sum_{j=1}^N m_j|\Sigma_j|,
  \qquad m_j\in\mathbb N,
\]
where the $\Sigma_j$ are smooth embedded closed minimal components.  By
\eqref{eq:willmore}, $\Area(\Sigma_j)\geq4\pi$ for every $j$.  Since the
total mass is exactly $4\pi$, there can be only one component, its
multiplicity is one, and equality holds in the Willmore inequality.  The
equality case makes $\Sigma_1$ a geodesic sphere.  Minimality then forces
it to be an equator.
\end{proof}

This argument also removes three possible ambiguities at once.  Allowing
integer multiplicity cannot produce another mass-$4\pi$ varifold;
allowing disconnected support cannot fit within the mass budget; and in
dimension three the min--max support has no singular set.  Even if one
reads the conjecture as asking for an arbitrary closed immersed minimal
surface of parametrized area $\omega_2$, rather than a particular min--max
varifold, the equality case of \eqref{eq:willmore} forces the immersion to
be an embedding onto a geodesic two-sphere.  Minimality makes it an equator;
equivalently, its induced varifold is $|E|$.

\section{Index and the lower-area obstruction}

Let $E\cong\Sph^2$ be an equator.  Its second fundamental form vanishes,
and $\Ric_{\Sph^3}(\nu,\nu)=2$.  With the convention that the Morse index
is the number of negative directions of the second-variation form,
\[
  Q_E(f,f)
  =\int_E\left(|\nabla f|^2-
  \bigl(|A_E|^2+\Ric(\nu,\nu)\bigr)f^2\right)d\mu
  =\int_{\Sph^2}\left(|\nabla f|^2-2f^2\right)d\mu.
\]
The spectrum of the nonnegative Laplacian $-\Delta_{\Sph^2}$ is
\[
  \ell(\ell+1),
  \qquad \ell=0,1,2,\ldots,
\]
with multiplicity $2\ell+1$.  Hence the constant mode $\ell=0$ is the
only negative direction, while the three $\ell=1$ modes span the kernel.
Therefore
\begin{equation}\label{eq:index}
  \ind(E)=1,
  \qquad
  \nul(E)=3.
\end{equation}

Combining Lemmas~\ref{lem:width} and \ref{lem:rigidity} with
\eqref{eq:index}, every minimal hypersurface realizing
$\omega_2(\Sph^3)=4\pi$ has index $1$.  In particular, no such realizer
has the index $p=2$ required by the conjecture.

There is a second, logically independent contradiction.  If
$\Sigma_+\subset\Sph^3$ were a nonempty closed minimal surface, then
\eqref{eq:willmore} would give
\[
  \Area(\Sigma_+)\geq4\pi=\Area(E).
\]
Thus no minimal hypersurface of strictly smaller area exists.  Here, as is
standard, ``hypersurface'' means a nonempty geometric hypersurface.  A
Brakke flow issuing from an equator may become extinct, but the empty
varifold is not the lower-area minimal hypersurface required by the
published statement.  This completes the proof of
Theorem~\ref{thm:main}.

\section{Scope of the correction}

The counterexample concerns the literal positive-Ricci alternative in
the published conjecture.  It does not disprove a version restricted to
generic bumpy metrics.  Nor does it settle the existence of connecting
flows when an appropriate higher-index width realizer and a lower-area
minimal endpoint are already known to exist.

In their Berger-sphere setting, Chen and Gaspar explain that approximation
and compactness produce a width realizer satisfying
\[
  \ind(\Sigma)\leq p\leq\ind(\Sigma)+\nul(\Sigma),
\]
rather than $\ind(\Sigma)=p$; see \cite[Remark~1]{ChenGaspar}.  The same
relation holds for the round equator at $p=2$, since $1\leq2\leq4$ by
\eqref{eq:index}.  Replacing the exact index condition by this
augmented-index relation in an appropriately qualified conjecture would
remove the first obstruction, but not the absence of a lower-area endpoint
on the bottom width plateau.  Any repaired statement must therefore also
address width plateaux or impose a condition ensuring that a lower-area
minimal critical point exists.

\begin{thebibliography}{9}

\bibitem{ChenGaspar}
J.~Chen and P.~Gaspar,
\emph{When can minimal hypersurfaces be connected by mean curvature flow?},
Math. Ann. \textbf{395} (2026), Article 40.
\href{https://doi.org/10.1007/s00208-026-03479-5}
{doi:10.1007/s00208-026-03479-5}.

\bibitem{MarquesNevesPositiveRicci}
F.~C. Marques and A.~Neves,
\emph{Existence of infinitely many minimal hypersurfaces in positive
Ricci curvature},
Invent. Math. \textbf{209} (2017), no.~2, 577--616.
\href{https://doi.org/10.1007/s00222-016-0716-3}
{doi:10.1007/s00222-016-0716-3}.

\bibitem{MarquesNevesWillmore}
F.~C. Marques and A.~Neves,
\emph{Min--max theory and the Willmore conjecture},
Ann. of Math. (2) \textbf{179} (2014), no.~2, 683--782.
\href{https://doi.org/10.4007/annals.2014.179.2.6}
{doi:10.4007/annals.2014.179.2.6}.

\bibitem{Nurser}
C.~A.~G. Nurser,
\emph{Low min--max widths of the round three-sphere},
Ph.D. thesis, Imperial College London, 2016.

\bibitem{Pitts}
J.~T. Pitts,
\emph{Existence and Regularity of Minimal Surfaces on Riemannian Manifolds},
Mathematical Notes, vol.~27, Princeton University Press, Princeton, 1981.

\bibitem{Willmore}
T.~J. Willmore,
\emph{Note on embedded surfaces},
An. \c{S}tiin\c{t}. Univ. ``Al. I. Cuza'' Ia\c{s}i Sec\c{t}. I a Mat.
\textbf{11B} (1965), 493--496.

\end{thebibliography}

\end{document}
